\iftestbuild % TEST-ONLY-BEGIN
\setcounter{section}{1}\setcounter{theorem}{23}
\fi % TEST-ONLY-END
% Source: book pp. 346--353 (PDF 360--367); no solutions.
% \exnote: Axler's small italic comment under an exercise
\section{Alternating Multilinear Forms}

\subsection{Multilinear Forms}

\begin{definition}[$V^m$]\label{ax:9.24}
For $m$ a positive integer, define $V^m$ by
\[ V^m = \ftimes{V\times\cdots\times V}{m\text{ times}}. \]
\end{definition}

Now we can define $m$-linear forms as a generalization of the bilinear forms
that we discussed in the previous section.

\begin{definition}[$m$-linear form, $V^{(m)}$, multilinear form]\label{ax:9.25}\leavevmode
\begin{itemize}
\item For $m$ a positive integer, an \emph{$m$-linear form} on $V$ is a
  function $\beta\colon V^m\to\F$ that is linear in each slot when the other
  slots are held fixed. This means that for each $k\in\{1,\dots,m\}$ and all
  $u_1,\dots,u_m\in V$, the function
  \[ v\mapsto\beta(u_1,\dots,u_{k-1},v,u_{k+1},\dots,u_m) \]
  is a linear map from $V$ to $\F$.
\item The set of $m$-linear forms on $V$ is denoted by $V^{(m)}$.
\item A function $\beta$ is called a \emph{multilinear form} on $V$ if it is an
  $m$-linear form on $V$ for some positive integer $m$.
\end{itemize}
\end{definition}

In the definition above, the expression
$\beta(u_1,\dots,u_{k-1},v,\allowbreak u_{k+1},\dots,u_m)$ means
$\beta(v,\allowbreak u_2,\dots,u_m)$
if $k=1$ and means $\beta(u_1,\dots,u_{m-1},v)$ if $k=m$.

A $1$-linear form on $V$ is a linear functional on $V$. A $2$-linear form on
$V$ is a bilinear form on $V$. You can verify that with the usual addition and
scalar multiplication of functions, $V^{(m)}$ is a vector space.

\begin{example}[$m$-linear forms]\label{ax:9.26}\leavevmode
\begin{itemize}
\item Suppose $\alpha,\rho\in V^{(2)}$. Define a function
  $\beta\colon V^4\to\F$ by
  \[ \beta(v_1,v_2,v_3,v_4) = \alpha(v_1,v_2)\,\rho(v_3,v_4). \]
  Then $\beta\in V^{(4)}$.
\item Define $\beta\colon\bigl(\Lin(V)\bigr)^m\to\F$ by
  \[ \beta(T_1,\dots,T_m) = \tr(T_1\cdots T_m). \]
  Then $\beta$ is an $m$-linear form on $\Lin(V)$.
\end{itemize}
\end{example}

Alternating multilinear forms, which we now define, play an important role as
we head toward defining determinants.

\begin{definition}[Alternating forms, $V_{\mathrm{alt}}^{(m)}$]\label{ax:9.27}
Suppose $m$ is a positive integer.
\begin{itemize}
\item An $m$-linear form $\alpha$ on $V$ is called \emph{alternating} if
  $\alpha(v_1,\dots,v_m)=0$ whenever $v_1,\dots,v_m$ is a list of vectors in
  $V$ with $v_j=v_k$ for some two distinct values of $j$ and $k$ in
  $\{1,\dots,m\}$.
\item $V_{\mathrm{alt}}^{(m)} = \{\alpha\in V^{(m)} : \alpha \text{ is an
  alternating } m\text{-linear form on } V\}$.
\end{itemize}
\end{definition}

You should verify that $V_{\mathrm{alt}}^{(m)}$ is a subspace of $V^{(m)}$.
See Example~\ref{ax:9.15} for examples of alternating $2$-linear forms. See
Exercise~2 for an example of an alternating $3$-linear form.

The next result tells us that if a linearly dependent list is input to an
alternating multilinear form, then the output equals $0$.

\begin{result}[Alternating multilinear forms and linear dependence]\label{ax:9.28}
Suppose $m$ is a positive integer and $\alpha$ is an alternating $m$-linear
form on $V$. If $v_1,\dots,v_m$ is a linearly dependent list in $V$, then
\[ \alpha(v_1,\dots,v_m) = 0. \]
\end{result}

\begin{proof}
Suppose $v_1,\dots,v_m$ is a linearly dependent list in $V$. By the linear
dependence lemma (\ref{ax:2.19}), some $v_k$ is a linear combination of
$v_1,\dots,v_{k-1}$. Thus there exist $b_1,\dots,b_{k-1}$ such that
$v_k=b_1v_1+\cdots+b_{k-1}v_{k-1}$. Now
\begin{align*}
\alpha(v_1,\dots,v_m)
  &= \alpha\biggl(v_1,\dots,v_{k-1},\sum_{j=1}^{k-1}b_jv_j,v_{k+1},\dots,v_m\biggr)\\
  &= \sum_{j=1}^{k-1}b_j\,\alpha(v_1,\dots,v_{k-1},v_j,v_{k+1},\dots,v_m)\\
  &= 0. \qedhere
\end{align*}
\end{proof}

The next result states that if $m>\dim V$, then there are no alternating
$m$-linear forms on $V$ other than the function on $V^m$ that is identically
$0$.

\begin{result}[No nonzero alternating $m$-linear forms for $m>\dim V$]\label{ax:9.29}
Suppose $m>\dim V$. Then $0$ is the only alternating $m$-linear form on $V$.
\end{result}

\begin{proof}
Suppose that $\alpha$ is an alternating $m$-linear form on $V$ and
$v_1,\dots,v_m\in V$. Because $m>\dim V$, this list is not linearly
independent (by \ref{ax:2.22}). Thus \ref{ax:9.28} implies that
$\alpha(v_1,\dots,v_m)=0$. Hence $\alpha$ is the zero function from $V^m$ to
$\F$.
\end{proof}

\subsection{Alternating Multilinear Forms and Permutations}

\begin{result}[Swapping input vectors in an alternating multilinear form]\label{ax:9.30}
Suppose $m$ is a positive integer, $\alpha$ is an alternating $m$-linear form
on $V$, and $v_1,\dots,v_m$ is a list of vectors in $V$. Then swapping the
vectors in any two slots of $\alpha(v_1,\dots,v_m)$ changes the value of
$\alpha$ by a factor of $-1$.
\end{result}

\begin{proof}
Put $v_1+v_2$ in both the first two slots, getting
\[ 0 = \alpha(v_1+v_2,v_1+v_2,v_3,\dots,v_m). \]
Use the multilinear properties of $\alpha$ to expand the right side of the
equation above (as in the proof of \ref{ax:9.16}) to get
\[ \alpha(v_2,v_1,v_3,\dots,v_m) = -\alpha(v_1,v_2,v_3,\dots,v_m). \]

Similarly, swapping the vectors in any two slots of $\alpha(v_1,\dots,v_m)$
changes the value of $\alpha$ by a factor of $-1$.
\end{proof}

To see what can happen with multiple swaps, suppose $\alpha$ is an alternating
$3$-linear form on $V$ and $v_1,v_2,v_3\in V$. To evaluate
$\alpha(v_3,v_1,v_2)$ in terms of $\alpha(v_1,v_2,v_3)$, start with
$\alpha(v_3,v_1,v_2)$ and swap the entries in the first and third slots,
getting
$\alpha(v_3,\allowbreak v_1,v_2)=-\alpha(v_2,v_1,v_3)$. Now in the last
expression, swap the entries in the first and second slots, getting
\[ \alpha(v_3,v_1,v_2) = -\alpha(v_2,v_1,v_3) = \alpha(v_1,v_2,v_3). \]
More generally, we see that if we do an odd number of swaps, then the value of
$\alpha$ changes by a factor of $-1$, and if we do an even number of swaps,
then the value of $\alpha$ does not change.

To deal with arbitrary multiple swaps, we need a bit of information about
permutations.

\begin{definition}[Permutation, $\operatorname{perm} m$]\label{ax:9.31}
Suppose $m$ is a positive integer.
\begin{itemize}
\item A \emph{permutation} of $(1,\dots,m)$ is a list $(j_1,\dots,j_m)$ that
  contains each of the numbers $1,\dots,m$ exactly once.
\item The set of all permutations of $(1,\dots,m)$ is denoted by
  $\operatorname{perm} m$.
\end{itemize}
\end{definition}

For example, $(2,3,4,5,1)\in\operatorname{perm} 5$. You should think of an
element of $\operatorname{perm} m$ as a rearrangement of the first $m$
positive integers.

The number of swaps used to change a permutation $(j_1,\dots,j_m)$ to the
standard order $(1,\dots,m)$ can depend on the specific swaps selected. The
following definition has the advantage of assigning a well-defined sign to
every permutation.

\begin{definition}[Sign of a permutation]\label{ax:9.32}
The \emph{sign} of a permutation $(j_1,\dots,j_m)$ is defined by
\[ \sgn(j_1,\dots,j_m) = (-1)^N, \]
where $N$ is the number of pairs of integers $(k,\ell)$ with
$1\le k<\ell\le m$ such that $k$ appears after $\ell$ in the list
$(j_1,\dots,j_m)$.
\end{definition}

Hence the sign of a permutation equals $1$ if the natural order has been
changed an even number of times and equals $-1$ if the natural order has been
changed an odd number of times.

\begin{example}[Signs]\label{ax:9.33}\leavevmode
\begin{itemize}
\item The permutation $(1,\dots,m)$ [no changes in the natural order] has
  sign $1$.
\item The only pair of integers $(k,\ell)$ with $k<\ell$ such that $k$
  appears after $\ell$ in the list $(2,1,3,4)$ is $(1,2)$. Thus the
  permutation $(2,1,3,4)$ has sign $-1$.
\item In the permutation $(2,3,\dots,m,1)$, the only pairs $(k,\ell)$ with
  $k<\ell$ that appear with changed order are $(1,2),(1,3),\dots,(1,m)$.
  Because we have $m-1$ such pairs, the sign of this permutation equals
  $(-1)^{m-1}$.
\end{itemize}
\end{example}

\begin{result}[Swapping two entries in a permutation]\label{ax:9.34}
Swapping two entries in a permutation multiplies the sign of the permutation
by $-1$.
\end{result}

\begin{proof}
Suppose we have two permutations, where the second permutation is obtained
from the first by swapping two entries. The two swapped entries were in their
natural order in the first permutation if and only if they are not in their
natural order in the second permutation. Thus we have a net change (so far)
of $1$ or $-1$ (both odd numbers) in the number of pairs not in their natural
order.

Consider each entry between the two swapped entries. If an intermediate entry
was originally in the natural order with respect to both swapped entries, then
it is now in the natural order with respect to neither swapped entry.
Similarly, if an intermediate entry was originally in the natural order with
respect to neither of the swapped entries, then it is now in the natural order
with respect to both swapped entries. If an intermediate entry was originally
in the natural order with respect to exactly one of the swapped entries, then
that is still true. Thus the net change (for each pair containing an entry
between the two swapped entries) in the number of pairs not in their natural
order is $2$, $-2$, or $0$ (all even numbers).

For all other pairs of entries, there is no change in whether or not they are
in their natural order. Thus the total net change in the number of pairs not
in their natural order is an odd number. Hence the sign of the second
permutation equals $-1$ times the sign of the first permutation.
\end{proof}

\begin{result}[Permutations and alternating multilinear forms]\label{ax:9.35}
Suppose $m$ is a positive integer and $\alpha\in V_{\mathrm{alt}}^{(m)}$. Then
\[ \alpha(v_{j_1},\dots,v_{j_m})
   = \bigl(\sgn(j_1,\dots,j_m)\bigr)\alpha(v_1,\dots,v_m) \]
for every list $v_1,\dots,v_m$ of vectors in $V$ and all
$(j_1,\dots,j_m)\in\operatorname{perm} m$.
\end{result}

\begin{proof}
Suppose $v_1,\dots,v_m\in V$ and $(j_1,\dots,j_m)\in\operatorname{perm} m$. We
can get from $(j_1,\dots,j_m)$ to $(1,\dots,m)$ by a series of swaps of entries
in different slots. Each such swap changes the value of $\alpha$ by a factor of
$-1$ (by \ref{ax:9.30}) and also changes the sign of the remaining permutation
by a factor of $-1$ (by \ref{ax:9.34}). After an appropriate number of swaps,
we reach the permutation $1,\dots,m$, which has sign $1$. Thus the value of
$\alpha$ changed signs an even number of times if $\sgn(j_1,\dots,j_m)=1$ and
an odd number of times if $\sgn(j_1,\dots,j_m)=-1$, which gives the desired
result.
\end{proof}

Our use of permutations now leads in a natural way to the following beautiful
formula for alternating $n$-linear forms on an $n$-dimensional vector space.

\begin{result}[Formula for $(\dim V)$-linear alternating forms on $V$]\label{ax:9.36}
Let $n=\dim V$. Suppose $e_1,\dots,\allowbreak e_n$ is a basis of $V$ and
$v_1,\dots,v_n\in V$. For each $k\in\{1,\dots,n\}$, let
$b_{1,k},\dots,b_{n,k}\in\F$ be such that
\[ v_k = \sum_{j=1}^n b_{j,k}e_j. \]
Then
\[ \alpha(v_1,\dots,v_n) = \alpha(e_1,\dots,e_n)
   \sum_{(j_1,\dots,j_n)\in\operatorname{perm} n}
   \bigl(\sgn(j_1,\dots,j_n)\bigr)b_{j_1,1}\cdots b_{j_n,n} \]
for every alternating $n$-linear form $\alpha$ on $V$.
\end{result}

\begin{proof}
Suppose $\alpha$ is an alternating $n$-linear form on $V$. Then
\begin{align*}
\alpha(v_1,\dots,v_n)
  &= \alpha\biggl(\sum_{j_1=1}^n b_{j_1,1}e_{j_1},\dots,
     \sum_{j_n=1}^n b_{j_n,n}e_{j_n}\biggr)\\
  &= \sum_{j_1=1}^n\cdots\sum_{j_n=1}^n b_{j_1,1}\cdots b_{j_n,n}\,
     \alpha(e_{j_1},\dots,e_{j_n})\\
  &= \sum_{(j_1,\dots,j_n)\in\operatorname{perm} n} b_{j_1,1}\cdots b_{j_n,n}\,
     \alpha(e_{j_1},\dots,e_{j_n})\\
  &= \alpha(e_1,\dots,e_n)\sum_{(j_1,\dots,j_n)\in\operatorname{perm} n}
     \bigl(\sgn(j_1,\dots,j_n)\bigr)b_{j_1,1}\cdots b_{j_n,n},
\end{align*}
where the third line holds because $\alpha(e_{j_1},\dots,e_{j_n})=0$ if
$j_1,\dots,j_n$ are not distinct integers, and the last line holds by
\ref{ax:9.35}.
\end{proof}

The following result will be the key to our definition of the determinant in
the next section.

\begin{result}[$\dim V_{\mathrm{alt}}^{(\dim V)}=1$]\label{ax:9.37}
The vector space $V_{\mathrm{alt}}^{(\dim V)}$ has dimension one.
\end{result}

\begin{proof}
Let $n=\dim V$. Suppose $\alpha$ and $\alpha'$ are alternating $n$-linear forms
on $V$ with $\alpha\ne0$. Let $e_1,\dots,e_n$ be such that
$\alpha(e_1,\dots,e_n)\ne0$. There exists $c\in\F$ such that
\[ \alpha'(e_1,\dots,e_n) = c\alpha(e_1,\dots,e_n). \]
Furthermore, \ref{ax:9.28} implies that $e_1,\dots,e_n$ is linearly
independent and thus is a basis of $V$.

Suppose $v_1,\dots,v_n\in V$. Let $b_{j,k}$ be as in \ref{ax:9.36} for
$j,k=1,\dots,n$. Then
\begin{align*}
\alpha'(v_1,\dots,v_n)
  &= \alpha'(e_1,\dots,e_n)\sum_{(j_1,\dots,j_n)\in\operatorname{perm} n}
     \bigl(\sgn(j_1,\dots,j_n)\bigr)b_{j_1,1}\cdots b_{j_n,n}\\
  &= c\alpha(e_1,\dots,e_n)\sum_{(j_1,\dots,j_n)\in\operatorname{perm} n}
     \bigl(\sgn(j_1,\dots,j_n)\bigr)b_{j_1,1}\cdots b_{j_n,n}\\
  &= c\alpha(v_1,\dots,v_n),
\end{align*}
where the first and last lines above come from \ref{ax:9.36}. The equation
above implies that $\alpha'=c\alpha$. Thus $\alpha',\alpha$ is not a linearly
independent list, which implies that $\dim V_{\mathrm{alt}}^{(n)}\le1$.

To complete the proof, we only need to show that there exists a nonzero
alternating $n$-linear form $\alpha$ on $V$ (thus eliminating the possibility
that $\dim V_{\mathrm{alt}}^{(n)}$ equals $0$). To do this, let
$e_1,\dots,e_n$ be any basis of $V$, and let $\varphi_1,\dots,\varphi_n\in V'$
be the linear functionals on $V$ that allow us to express each element of $V$
as a linear combination of $e_1,\dots,e_n$. In other words,
\[ v = \sum_{j=1}^n\varphi_j(v)e_j \]
for every $v\in V$ (see \ref{ax:3.114}). Now for $v_1,\dots,v_n\in V$, define
\refstepcounter{theorem}%
\[ \alpha(v_1,\dots,v_n) = \sum_{(j_1,\dots,j_n)\in\operatorname{perm} n}
   \bigl(\sgn(j_1,\dots,j_n)\bigr)\varphi_{j_1}(v_1)\cdots\varphi_{j_n}(v_n).
   \tag*{\thetheorem}\label{ax:9.38} \]
The verification that $\alpha$ is an $n$-linear form on $V$ is
straightforward.

To see that $\alpha$ is alternating, suppose $v_1,\dots,v_n\in V$ with
$v_1=v_2$. For each $(j_1,\dots,j_n)\in\operatorname{perm} n$, the permutation
$(j_2,j_1,j_3,\dots,j_n)$ has the opposite sign. Because $v_1=v_2$, the
contributions from these two permutations to the sum in \ref{ax:9.38} cancel
each other. Hence $\alpha(v_1,v_1,v_3,\dots,v_n)=0$. Similarly,
$\alpha(v_1,\dots,v_n)=0$ if any two vectors in the list $v_1,\dots,v_n$ are
equal. Thus $\alpha$ is alternating.

Finally, consider \ref{ax:9.38} with each $v_k=e_k$. Because
$\varphi_j(e_k)$ equals $0$ if $j\ne k$ and equals $1$ if $j=k$, only the
permutation $(1,\dots,n)$ makes a nonzero contribution to the right side of
\ref{ax:9.38} in this case, giving the equation $\alpha(e_1,\dots,e_n)=1$.
Thus we have produced a nonzero alternating $n$-linear form $\alpha$ on $V$,
as desired.
\end{proof}

\marginnote{The formula \ref{ax:9.38} used in the last proof to construct a
nonzero alternating $n$-linear form came from the formula in \ref{ax:9.36},
and that formula arose naturally from the properties of an alternating
multilinear form.}%
Earlier we showed that the value of an alternating multilinear form applied to
a linearly dependent list is $0$; see \ref{ax:9.28}. The next result provides a
converse of \ref{ax:9.28} for alternating $n$-linear forms when $n=\dim V$. In
the following result, the statement that $\alpha$ is nonzero means (as usual
for a function) that $\alpha$ is not the function on $V^n$ that is
identically $0$.

\begin{result}[Alternating $(\dim V)$-linear forms and linear independence]\label{ax:9.39}
Let $n=\dim V$. Suppose $\alpha$ is a nonzero alternating $n$-linear form on
$V$ and $e_1,\dots,e_n$ is a list of vectors in $V$. Then
\[ \alpha(e_1,\dots,e_n)\ne0 \]
if and only if $e_1,\dots,e_n$ is linearly independent.
\end{result}

\begin{proof}
First suppose $\alpha(e_1,\dots,e_n)\ne0$. Then \ref{ax:9.28} implies that
$e_1,\dots,e_n$ is linearly independent.

To prove the implication in the other direction, now suppose $e_1,\dots,e_n$
is linearly independent. Because $n=\dim V$, this implies that
$e_1,\dots,e_n$ is a basis of $V$ (see \ref{ax:2.38}).

Because $\alpha$ is not the zero $n$-linear form, there exist
$v_1,\dots,v_n\in V$ such that $\alpha(v_1,\dots,v_n)\ne0$. Now \ref{ax:9.36}
implies that $\alpha(e_1,\dots,e_n)\ne0$.
\end{proof}

% ------------------------------------------------------------------ exercises
\exercisesheading{9B}

\begin{exercise}
Suppose $m$ is a positive integer. Show that $\dim V^{(m)}=(\dim V)^m$.
\end{exercise}

\begin{exercise}
Suppose $n\ge3$ and $\alpha\colon\F^n\times\F^n\times\F^n\to\F$ is defined by
\begin{multline*}
\alpha\bigl((x_1,\dots,x_n),(y_1,\dots,y_n),(z_1,\dots,z_n)\bigr)\\
  = x_1y_2z_3-x_2y_1z_3-x_3y_2z_1-x_1y_3z_2+x_3y_1z_2+x_2y_3z_1.
\end{multline*}
Show that $\alpha$ is an alternating $3$-linear form on $\F^n$.
\end{exercise}

\begin{exercise}
Suppose $m$ is a positive integer and $\alpha$ is an $m$-linear form on $V$
such that $\alpha(v_1,\dots,v_m)=0$ whenever $v_1,\dots,v_m$ is a list of
vectors in $V$ with $v_j=v_{j+1}$ for some $j\in\{1,\dots,m-1\}$. Prove that
$\alpha$ is an alternating $m$-linear form on $V$.
\end{exercise}

\begin{exercise}
Prove or give a counterexample: If $\alpha\in V_{\mathrm{alt}}^{(4)}$, then
\[ \{(v_1,v_2,v_3,v_4)\in V^4 : \alpha(v_1,v_2,v_3,v_4)=0\} \]
is a subspace of $V^4$.
\end{exercise}

\begin{exercise}
Suppose $m$ is a positive integer and $\beta$ is an $m$-linear form on $V$.
Define an $m$-linear form $\alpha$ on $V$ by
\[ \alpha(v_1,\dots,v_m) = \sum_{(j_1,\dots,j_m)\in\operatorname{perm} m}
   \bigl(\sgn(j_1,\dots,j_m)\bigr)\beta(v_{j_1},\dots,v_{j_m}) \]
for $v_1,\dots,v_m\in V$. Explain why $\alpha\in V_{\mathrm{alt}}^{(m)}$.
\end{exercise}

\begin{exercise}
Suppose $m$ is a positive integer and $\beta$ is an $m$-linear form on $V$.
Define an $m$-linear form $\alpha$ on $V$ by
\[ \alpha(v_1,\dots,v_m) = \sum_{(j_1,\dots,j_m)\in\operatorname{perm} m}
   \beta(v_{j_1},\dots,v_{j_m}) \]
for $v_1,\dots,v_m\in V$. Explain why
\[ \alpha(v_{k_1},\dots,v_{k_m}) = \alpha(v_1,\dots,v_m) \]
for all $v_1,\dots,v_m\in V$ and all $(k_1,\dots,k_m)\in\operatorname{perm} m$.
\end{exercise}

\begin{exercise}
Give an example of a nonzero alternating $2$-linear form $\alpha$ on $\R^3$
and a linearly independent list $v_1,v_2$ in $\R^3$ such that
$\alpha(v_1,v_2)=0$.
\exnote{This exercise shows that \ref{ax:9.39} can fail if the hypothesis that
$n=\dim V$ is deleted.}
\end{exercise}
